\cleardoublepage % memastikan bab baru dimulai di halaman ganjil (kanan)
\chapter{PENDAHULUAN}
\label{chap:pendahuluan}

% --- Latar Belakang ---
\section{Latar Belakang}
% Hipertensi atau tekanan darah tinggi merupakan masalah kesehatan global yang menjadi faktor risiko utama untuk berbagai penyakit kardiovaskular, termasuk stroke dan gagal jantung \autocite{Mills2020Global}. Menurut data Organisasi Kesehatan Dunia, diperkirakan 1,28 miliar orang dewasa di seluruh dunia menderita hipertensi \autocite{WHO2023}. Di Indonesia, data Kementerian Kesehatan RI juga mencatat prevalensi yang signifikan \autocite{ri2018}. Meskipun manajemen hipertensi telah menjadi prioritas, penanganan yang umum terjadi saat ini sering kali bersifat reaktif \autocite{Whelton2018}. Intervensi medis baru dilakukan setelah pasien mengalami gejala atau setelah Tekanan Darah (TD) terdeteksi sangat tinggi, yang berpotensi berujung pada komplikasi akut seperti krisis hipertensi \textcite{Haddadin2018Hypertensive}.
Hipertensi atau tekanan darah tinggi merupakan masalah kesehatan global yang menjadi faktor risiko utama untuk berbagai penyakit kardiovaskular, termasuk stroke dan gagal jantung \autocite{Mills2020Global}. Selain itu, berdasarkan estimasi dari \textcite{WHO2023}, penderita hipertensi di seluruh dunia telah mencapai angka 1,28 miliar orang pada kelompok usia dewasa. \textcite{ri2018} juga mencatat bahwa angka kejadian hipertensi di dalam negeri telah menunjukkan prevalensi yang tinggi dan meningkat secara signifikan dari waktu ke waktu. Meskipun penanganan hipertensi telah menjadi prioritas, praktik yang sering dilakukan saat ini masih cenderung bersifat reaktif, yaitu baru dilakukan setelah masalah atau komplikasi muncul \autocite{Whelton2018}. Hal tersebut selaras dengan kutipan oleh \textcite{Haddadin2018Hypertensive} yang menyampaikan bahwa intervensi medis biasanya baru berjalan setelah kondisi pasien memburuk secara klinis atau tekanan darahnya terdeteksi sangat tinggi. Pola penanganan yang reaktif ini berisiko besar menyeret pasien ke dalam kondisi darurat seperti krisis hipertensi dan komplikasi lainnya.
% Pendekatan berbasis \textit{machine learning} menawarkan alternatif yang menjanjikan dalam memprediksi risiko hipertensi. Algoritma ML terkini memiliki kemampuan superior dalam memodelkan hubungan non-linier yang kompleks antara berbagai variabel risiko \autocite{Tarimana2024}. Studi terbaru menunjukkan bahwa model ML mampu mencapai akurasi klasifikasi risiko yang jauh lebih tinggi dibandingkan regresi logistik tradisional. Namun, penerapan model-model ini di ranah klinis menghadapi kendala besar yaitu fenomena \textit{black box} yang membuat model menghasilkan prediksi tanpa menyertakan alasan yang dapat dipahami manusia.
Oleh karena itu, pendekatan berbasis \textit{machine learning} kini menjadi salah satu  opsi yang sangat potensial dalam melakukan prediksi terhadap risiko hipertensi. Algoritma ML modern terbukti sangat andal dalam memetakan interaksi non-linier yang rumit di antara beragam variabel risiko yang ada \autocite{Tarimana2024}. Namun, implementasi praktis model tersebut di lingkungan klinis masih terkendala oleh sifat black box, dengan sistem masih mungkin untuk mengeluarkan hasil prediksi tanpa disertai dasar penalaran yang logis bagi penggunanya. menurut \textcite{amann2022explain}, istilah \textit{black box}


Dalam konteks medis, hasil prediksi saja tidak cukup. Tenaga medis dan pasien membutuhkan interpretabilitas untuk memahami alasan di balik sebuah prediksi agar dapat mempercayai dan menindaklanjuti hasil tersebut. Kompleksitas algoritma prediksi yang bersifat \textit{black box} sering kali menjadi penghambat adopsi teknologi ini dalam praktik klinis sehari-hari karena kurangnya transparansi dan kepercayaan dari pengguna \autocite{Adadi2018}.

Sebagai solusi atas kendala tersebut, paradigma Explainable Artificial Intelligence atau XAI diterapkan untuk membuka mekanisme internal model sehingga alur pengambilan keputusan dapat ditelusuri kembali secara logis \autocite{Adadi2018}. Transparansi ini merupakan elemen fundamental dalam informatika medis yang bertujuan menjamin akuntabilitas keputusan serta memastikan keselamatan pasien tetap menjadi prioritas utama \autocite{Holzinger2019}.

Meskipun metode XAI mampu memberikan transparansi teknis, luaran yang dihasilkan umumnya masih berupa data numerik atau visualisasi fitur yang kompleks bagi pengguna non-teknis. Guna menjembatani kesenjangan pemahaman tersebut, teknologi \textit{Large Language Model} (LLM) dimanfaatkan dalam penelitian ini untuk mentransformasi hasil interpretasi model menjadi narasi bahasa alami yang personal \autocite{Thirunavukarasu2023Large}. Akan tetapi, tantangan utama dalam penerapan model bahasa generik pada domain sensitif adalah kecenderungan model untuk menghasilkan fakta atau halusinasi tidak berdasar \autocite{Ji2023Survey}. Sebagai strategi mitigasi yang efektif, arsitektur Retrieval-Augmented Generation (RAG) diterapkan untuk memaksa model melakukan validasi silang terhadap basis pengetahuan eksternal yang terpercaya sebelum menyusun jawaban akhir \autocite{Lewis2020Retrieval}.

Oleh karena itu, penelitian ini mengusulkan pendekatan baru dengan mengintegrasikan LLM berbasis arsitektur RAG ke dalam sistem prediksi hipertensi. LLM akan berfungsi untuk menjembatani penerjemahan luaran yang matematis dari XAI menjadi narasi bahasa alami yang personal dan mudah dipahami. Penggunaan RAG secara spesifik bertujuan untuk memitigasi risiko halusinasi pada LLM dengan membatasi respons model pada referensi pedoman medis yang terverifikasi. Berbeda dari penelitian sebelumnya yang umumnya hanya berfokus pada peningkatan performa prediktif atau penyajian interpretasi teknis, penelitian ini mengusulkan pendekatan terintegrasi yang mengombinasikan XAI dan LLM berbasis RAG untuk menghasilkan penjelasan klinis yang tervalidasi secara kontekstual, khususnya dalam konteks prediksi risiko hipertensi.

% --- Rumusan Masalah ---
\section{Rumusan Masalah}
Berdasarkan latar belakang yang telah diuraikan, maka rumusan masalah dalam penelitian ini adalah:
\begin{enumerate}
\item Bagaimana membangun model Machine Learning untuk prediksi risiko hipertensi yang mampu mendapatkan performa klasifikasi yang optimal berdasarkan dataset tabular dengan variabel risiko yang beragam?
\item Bagaimana menerapkan metode Explainable Artificial Intelligence (XAI) untuk membuka mekanisme \textit{black box} model terpilih guna menghasilkan gambaran fitur yang mempengaruhi prediksi dengan akurat?
\item Bagaimana menerapkan arsitektur RAG pada LLM untuk mentransformasi luaran teknis model menjadi penjelasan naratif yang personal dan tervalidasi oleh pedoman medis?
\item Bagaimana mengukur efektivitas keseluruhan sistem yang diusulkan jika dievaluasi berdasarkan gabungan metrik performa prediksi ML, kualitas interpretabilitas XAI, tingkat faktualitas LLM, dan relevansi narasi yang dihasilkan?
\end{enumerate}

% --- Tujuan ---
\section{Tujuan}
Tujuan utama pelaksanaan Tugas Akhir ini adalah:
\begin{enumerate}
\item Membangun dan mengoptimasi model klasifikasi risiko hipertensi menggunakan algoritma ML dengan target metrik evaluasi melampaui 70\%.
\item Menerapkan metode Explainable Artificial Intelligence (XAI) untuk menghasilkan atribusi fitur yang transparan dan dapat diinterpretasikan secara klinis.
\item Mengembangkan modul narasi berbasis LLM dengan arsitektur RAG yang memanfaatkan hasil XAI untuk memberikan penjelasan personal yang tervalidasi.
\end{enumerate}

% --- Batasan Masalah ---
\section{Batasan Masalah}
Penelitian ini menetapkan batasan-batasan spesifik untuk menjaga fokus pembahasan dan kelayakan pengerjaan dalam kerangka waktu yang tersedia. Batasan masalah dalam penelitian ini adalah sebagai berikut:
\begin{enumerate}
\item Penelitian ini menggunakan data sekunder terstruktur dalam bentuk tabular yang diperoleh dari repositori publik terverifikasi atau himpunan data klinis anonim. Variabel yang digunakan terbatas pada fitur demografis, riwayat medis, dan tanda vital klinis yang relevan dengan faktor risiko hipertensi.
\item Model prediksi dikembangkan menggunakan pendekatan Supervised Machine Learning dengan fokus pada kasus klasifikasi biner, yaitu membedakan antara kelas risiko hipertensi dan non-hipertensi. Evaluasi model difokuskan pada metrik performa teknis seperti akurasi, presisi, recall, dan F1-score.
\item Implementasi XAI dibatasi pada metode post-hoc model-agnostik, khususnya Shapley Additive Explanations (SHAP) atau Local Interpretable Model-agnostic Explanations (LIME), untuk menghasilkan nilai atribusi fitur. Salah satu dari metode ini dipilih untuk menjelaskan kontribusi setiap variabel terhadap prediksi akhir model tanpa mengubah struktur internal model prediksi itu sendiri.
\item Penelitian ini memanfaatkan LLM yang telah dilatih sebelumnya melalui antarmuka pemrograman aplikasi (API) atau model lokal yang sudah terkuantisasi. Penelitian tidak mencakup proses fine-tuning atau pelatihan ulang model bahasa secara keseluruhan, melainkan berfokus pada teknik prompt engineering dan integrasi konteks eksternal.
\item Sistem RAG menggunakan basis pengetahuan yang bersumber dari dokumen pedoman klinis resmi, seperti panduan dari Kementerian Kesehatan Republik Indonesia, World Health Organization (WHO), atau asosiasi kardiologi internasional yang relevan. Sistem tidak melakukan pencarian informasi secara langsung ke internet terbuka untuk menjamin validitas referensi.
\item Hasil akhir pengembangan sistem ini adalah purwarupa perangkat lunak berbasis web yang berfungsi sebagai alat skrining awal dan edukasi bagi masyarakat awam. Luaran dari sistem ini sama sekali tidak diperuntukkan sebagai diagnosis medis yang mutlak, melainkan hanya sebagai opini kedua, sehingga tidak dapat menggantikan diagnosis resmi dari dokter atau tenaga kesehatan profesional.
\end{enumerate}

% --- Metodologi Pengerjaan TA ---
\section{Metodologi}
Penelitian ini dilaksanakan dengan mengadopsi kerangka kerja standar CRISP-DM (Cross-Industry Standard Process for Data Mining). Kerangka kerja ini dipilih karena menyediakan struktur yang komprehensif dan iteratif, memastikan keselarasan antara tujuan klinis dengan teknis pengembangan model. Tahapan pelaksanaan penelitian dibagi menjadi enam tahap utama sebagai berikut:
\begin{enumerate}
\item \textbf{Business Understanding}\\
Tahap ini merupakan tahap fondasi yang bertujuan untuk mendefinisikan masalah penelitian secara presisi serta membangun landasan teoretis yang kuat. Pertama-tama, investigasi awal dilakukan secara mendalam untuk mengumpulkan fakta empiris yang mendasari latar belakang masalah. Kegiatan ini mencakup penelusuran data statistik kesehatan global dari World Health Organization dan data nasional dari Kementerian Kesehatan Republik Indonesia guna memvalidasi urgensi hipertensi sebagai komorbiditas utama. Selain itu, dilakukan analisis kesenjangan teknologi untuk mengidentifikasi hambatan adopsi kecerdasan artifisial di ranah klinis, khususnya terkait isu transparansi model dan risiko halusinasi, yang kemudian menjadi dasar perumusan solusi. Guna mendapatkan solusi terkini, dilakukan pencarian literatur pada basis data bereputasi seperti IEEE Xplore, PubMed, dan ScienceDirect. Pencarian informasi difokuskan pada pengumpulan teori dan metode state-of-the-art dengan menggunakan kata kunci spesifik pada komponen-komponen utama penelitian, yaitu hipertensi, ML, XAI, LLM, dan RAG. Literatur yang terkumpul kemudian melalui proses sintesis, yang hasilnya didokumentasikan secara rinci dalam Bab II Studi Literatur.

\item \textbf{Data Understanding}\\
Setelah tujuan penelitian ditetapkan, tahap selanjutnya berfokus pada akuisisi dan eksplorasi dataset. Penelitian ini menggunakan dataset sekunder yang diperoleh dari repositori publik terverifikasi. Setelah dataset sudah diakuisisi, dilakukan eksplorasi awal dan penyesuaian dataset agar dapat digunakan pada tahap pelatihan model. Apabila dataset yang tersedia belum mencukupi kebutuhan desain, maka dilakukan penyesuaian atau pengolahan tambahan.

\item \textbf{Data Preparation}\\
Tahap persiapan data bertujuan untuk mentransformasi data mentah menjadi format yang siap untuk proses pemodelan. Kegiatan teknis yang dilakukan meliputi \textit{data cleaning} untuk menangani anomali yang ditemukan pada tahap sebelumnya, serta pengkodean fitur kategorikal agar mengubah data non-numerik menjadi format numerik yang dapat dipahami dan digunakan oleh algoritma machine learning. Jika ditemukan ketidakseimbangan kelas pada label target, teknik penyeimbangan data akan diterapkan. Tahap ini diakhiri dengan pembagian dataset menjadi himpunan latih (training set) dan himpunan uji (test set) untuk digunakan model.

\item \textbf{Modelling}\\
Tahap ini merupakan inti dari pengembangan teknis yang terdiri dari tiga komponen utama yang saling terintegrasi. Pertama, dilakukan pelatihan menggunakan beberapa algoritma machine learning. Setelahnya, algoritma machine learning dibangun dan dilatih menggunakan dataset yang telah dipersiapkan. Performa setiap model dievaluasi menggunakan metrik ROC-AUC serta sejumlah metrik evaluasi lainnya. Model dengan hasil kinerja paling optimal kemudian dipilih untuk tahap interpretasi lebih lanjut. Kedua, metode XAI diimplementasikan pada model tersebut untuk menghasilkan nilai atribusi fitur yang menjelaskan alasan di balik setiap prediksi. Ketiga, pipeline Retrieval-Augmented Generation dikembangkan untuk menghubungkan luaran SHAP dengan Large Language Model. Sistem ini dikonfigurasi secara khusus agar merujuk pada basis pengetahuan pedoman medis sebelum menyusun narasi penjelasan, sehingga luaran yang dihasilkan tetap berada dalam koridor medis yang valid.

\item \textbf{Evaluation}\\
Tahap evaluasi dilakukan untuk mengukur efektivitas sistem secara menyeluruh sebelum tahap penyebaran. Evaluasi kuantitatif dilakukan terhadap model prediksi menggunakan metrik standar klasifikasi seperti akurasi, presisi, recall, F1-score, dan metrik lainnya. Sementara itu, evaluasi kualitatif difokuskan pada hasil interpretasi yang diperoleh dari metode XAI dengan melakukan validasi domain expert guna menilai tingkat kesesuaiannya dengan pengetahuan serta praktik di bidang terkait. Berdasarkan proses evaluasi tersebut, ditetapkan metode interpretasi yang paling relevan dan layak untuk diimplementasikan dalam sistem. Setelah itu penilaian kualitas narasi yang dihasilkan oleh Large Language Model. Aspek yang dinilai mencakup faktualitas medis, koherensi bahasa, dan relevansi terhadap konteks klinis pasien, dengan cara membandingkan luaran sistem terhadap referensi standar untuk memastikan tidak terjadi halusinasi informasi.

\item \textbf{Deployment}\\
Tahap terakhir dalam penelitian ini adalah pengembangan purwarupa perangkat lunak (software prototype) berbasis web. Model yang telah dilatih dan divalidasi diintegrasikan ke dalam antarmuka aplikasi yang ramah pengguna.
\end{enumerate}

% --- Sistematika Penulisan ---
\section{Sistematika Penulisan}
Sistematika penulisan laporan tugas akhir ini memberikan panduan bagi pembaca tentang struktur dan alur pembahasan. Berikut adalah sistematika penulisannya:
\begin{enumerate}
\item   Bab I Pendahuluan: Berisi latar belakang, rumusan masalah, tujuan, batasan masalah, metodologi, dan sistematika penulisan laporan tugas akhir.
\item   Bab II Studi Literatur: Berisi kajian pustaka yang relevan dengan topik tugas akhir, termasuk teori-teori, konsep-konsep, dan hasil-hasil penelitian sebelumnya yang mendukung pelaksanaan tugas akhir.
\item   Bab III Analisis: Berisi penjelasan tentang analisis masalah dan kebutuhan pengguna serta berbagai alternatif solusinya.
\item   Bab IV Perancangan: Berisi penjelasan tentang perancangan solusi yang diusulkan, misalnya arsitektur sistem, desain basis data, antarmuka pengguna, dan sebagainya.
\item   Bab V Implementasi: Berisi penjelasan tentang proses implementasi solusi yang diusulkan, termasuk langkah-langkah yang diambil, teknologi yang digunakan, dan hasil yang diperoleh.
\item   Bab VI Evaluasi: Berisi analisis dan evaluasi terhadap solusi yang diimplementasikan, termasuk metode pengujian, hal-hal yang disiapkan sebelum pengujian, hasil pengujian, dan penilaian kinerja.
\item   Bab VII Kesimpulan dan Saran: Berisi kesimpulan dari hasil pelaksanaan tugas akhir dan saran untuk pengembangan lebih lanjut.
\item   Daftar Pustaka: Berisi daftar referensi yang digunakan dalam penyusunan laporan tugas akhir.
\item   Lampiran: Berisi dokumen-dokumen pendukung seperti kode program, data lengkap hasil pengujian, rekaman wawancara, dan lain-lain. 
\end{enumerate}